\documentclass[dvipdfmx,usejistfm]{ieej-tec2}

% ============================================================
% Preamble
% ============================================================
\usepackage{preamble}
\usepackage{siunitx}
\usepackage{tabularx}

\graphicspath{{fig/}}
\DeclareSIUnit{\degr}{deg}

% ============================================================
% Paper information
% ============================================================
% 論文番号（1頁目右上）
\articlenumber{IIC-XX-XXX}

% タイトル
\jtitle{和文タイトル}
\etitle{English Title}

% 著者
\authorlist{%
  \authorentry{堀越太智*}{Taichi Horikoshi}{東京大学}%
  \authorentry{}{}{}%
}

% 所属
\affiliate[1]
 {所属機関名}
 {Affiliation in English}
% \affiliate[2]
%  {所属機関名2}
%  {Affiliation 2 in English}

% 必要に応じて所属表示位置を調整
\breakaffiliate{1}
% \breakaffiliate{2}

\begin{document}

% ============================================================
% Keywords and abstract
% ============================================================
\begin{jkeyword}
キーワード1,キーワード2,キーワード3
\end{jkeyword}

\begin{ekeyword}
keyword 1, keyword 2, keyword 3
\end{ekeyword}

\begin{abstract}
\ \ 英文アブストラクトを記述する。
\end{abstract}

\maketitle

% ============================================================
% Main text
% ============================================================
\section{序論}
研究背景，従来研究，課題，本稿の目的および構成を記述する。

% 引用例：\cite{ReferenceKey}
% 複数引用例：\cite{ReferenceKey1,ReferenceKey2}

\section{対象システム}
対象とするシステム，実験装置，モデル化，パラメータなどを記述する。

\subsection{実験装置の概要}
実験装置の構成を記述する。

% ------------------------------------------------------------
% Figure template
% ------------------------------------------------------------
% \begin{figure}[tb]
%   \centering
%   \includegraphics[width=0.7\linewidth]{figure_name.pdf}
%   \bicaption{和文キャプション}{English caption.}
%   \label{fig:figure_label}
% \end{figure}
%
% 本文中の参照例：図\ref{fig:figure_label}に示す。

% ------------------------------------------------------------
% Two-panel figure template
% ------------------------------------------------------------
% \begin{figure}[tb]
%   \centering
%   \begin{subcaptionblock}[t]{0.48\linewidth}
%     \centering
%     \includegraphics[width=\linewidth]{figure_a.pdf}
%     \bisubcaption{和文キャプションA}{English caption A}
%   \end{subcaptionblock}
%   \hfill
%   \begin{subcaptionblock}[t]{0.48\linewidth}
%     \centering
%     \includegraphics[width=\linewidth]{figure_b.pdf}
%     \bisubcaption{和文キャプションB}{English caption B}
%   \end{subcaptionblock}
%   \bicaption{全体の和文キャプション}{Overall English caption.}
%   \label{fig:two_panel_label}
% \end{figure}

\subsection{数理モデル}
対象の数理モデルを記述する。

% ------------------------------------------------------------
% Equation templates
% ------------------------------------------------------------
\begin{equation}
  y(s) = G(s)u(s)
  \label{eq:example}
\end{equation}

% 本文中の参照例：式(\ref{eq:example})より，以下が得られる。

% \begin{align}
%   x_1 &= a+b, \\
%   x_2 &= c+d
%   \label{eq:align_example}
% \end{align}

% \begin{equation}
%   A =
%   \begin{pmatrix}
%     a_{11} & a_{12} \\
%     a_{21} & a_{22}
%   \end{pmatrix}
%   \label{eq:matrix_example}
% \end{equation}

% ------------------------------------------------------------
% Table template
% ------------------------------------------------------------
% \begin{table}[tb]
%   \centering
%   \bicaption{和文表題}{English table caption.}
%   \begin{tabularx}{\linewidth}{Xll}
%     \hline
%     Parameter & Value & Unit \\
%     \hline
%     Parameter 1 & $0.00$ & \si{unit} \\
%     Parameter 2 & $0.00$ & \si{unit} \\
%     \hline
%   \end{tabularx}
%   \label{tab:parameter_label}
% \end{table}
%
% 本文中の参照例：表\ref{tab:parameter_label}に示す。

\section{提案手法}
提案手法の概要，理論，設計手順などを記述する。

\subsection{手法の概要}
提案手法の基本的な考え方を記述する。

\subsection{設計方法}
設計式やアルゴリズムを記述する。

\section{検証}
シミュレーションおよび実験による検証方法を記述する。

\subsection{検証条件}
シミュレーション条件，実験条件，評価指標などを記述する。

\subsection{シミュレーション結果}
シミュレーション結果と考察を記述する。

\subsection{実験結果}
実験結果と考察を記述する。

\section{結論}
本稿で得られた成果をまとめ，今後の課題を記述する。

% ============================================================
% Bibliography
% ============================================================
\bibliographystyle{IEEJtran}
\bibliography{cite}

\end{document}
